%%%%%%%%%%%%%%%%%%%%%%%%%%%%%%%%%%%%%%%%%%%%%%%%%%%%%%%%%%%%%%%%%%%%%%%%%%%%
% forced-objects.tex  (v4: rewritten for readability)
%
% One idea per paragraph, plain words first, symbols second.  Formal statements
% live in Boxes and in the Methods; the narrative carries the argument.
%
% HONESTY LEDGER (unchanged from v3):
%   * "checked" = Lean 4 declaration, no `sorry`, no named axiom beyond O1;
%     `native_decide` occurrences are listed in Methods and support results not
%     reported here.
%   * one named axiom, O1, and no theorem consumes it.
%   * precondition C0 (the framework is the standard one) is NOT formalised;
%     statements that need it are marked "modulo C0".
%   * the novelty measurement is refuted by the paper's own examples and is
%     reported as a failure, not as a working gate.
%%%%%%%%%%%%%%%%%%%%%%%%%%%%%%%%%%%%%%%%%%%%%%%%%%%%%%%%%%%%%%%%%%%%%%%%%%%%

\documentclass[aps,prd,preprint,superscriptaddress,amsmath,amssymb]{revtex4-2}
\emergencystretch=2.5em
\usepackage{graphicx}
\usepackage{bm}
\usepackage{dcolumn}
\usepackage{amsthm}
\usepackage{booktabs}
\usepackage{color}

\newtheorem{theorem}{Theorem}[section]
\newtheorem{definition}[theorem]{Definition}
\newtheorem{remark}[theorem]{Remark}
\newtheorem{obligation}[theorem]{Obligation}

\newcommand{\lean}[1]{\texttt{\small\detokenize{#1}}}

\begin{document}

\title{Where the Ladder Must Stand}

\author{Yuan-Jie Liu}
\email{yuanjieliu65@gmail.com}

\date{\today}

\begin{abstract}
Mathematicians often say that an object \emph{had to} exist: you cannot subtract
in the natural numbers, so the integers appear; you cannot solve $x^2+1=0$ in the
reals, so the complex numbers appear.  The same objects are also called
inventions, and both descriptions are usually offered without a test.  This paper
proposes one, in the form of three questions that can be answered yes or no, and
reports what happened when the test was turned on the author's own work.
Question~1: can the demand be met on the same elements?  Question~2: among all
extensions, is there a smallest one, determined up to a unique relabelling?
Question~3: is the result new, judged against a catalogue of known objects?  If
Question~2 is answered yes, the object is pinned down: any two smallest
extensions are the same up to a relabelling, and the relabelling is unique.
Three steps of the number tower pass Questions~1 and~2.  The step from the
rationals to the reals does not: its missing minimality is recorded as a named
obligation rather than proved, and no theorem depends on it.  A counting
invariant proposed inside the same programme is refused, because it collapses to
a difference of prime-factor counts and is therefore already definable on the
same elements; the refusal is argued by hand, since that layer contains no
formal demand.  Question~3 broke: the implementation of ``is it new?'' compared
the sizes of carriers, and the integers, the rationals and the complex numbers
all have the same size, so the test refused everything it was meant to admit.
The failure is reported at the same length as the successes, because it is the
clearest evidence that the test has teeth.  What survives is two working
questions, one specified but unimplemented third, and a short list of what this
approach cannot do, including the fact that every verdict is relative to a
stated triple of ambient theory, refinement and catalogue.  The mathematical
core is formalised in Lean~4 with Mathlib: no \texttt{sorry}, one named axiom,
and no theorem that consumes it.
\end{abstract}

\maketitle

\begin{quote}
\itshape
``Whatever can be defined may be questioned.''
\par
\vspace{0.4em}
--- \emph{Galactic Empire} (Isaac Asimov)
\end{quote}

\section{The question}
\label{sec:question}

\paragraph{A story with a wrong answer.}
Suppose a program announces that it has discovered a new number: $\alpha$, with
$\alpha^2=7$.  Nothing about this is wrong, and nothing about it is a discovery.
The rationals do not contain a square root of seven, so $\alpha$ does extend them;
but the extension is the field everyone already knows, with $\alpha$ written
instead of $\sqrt 7$.  Any test that only asks ``was there a solution before?''
will certify this as new, and it will certify a great deal else besides.

The same weakness appears in the standard justifications of the number systems.
``The natural numbers cannot solve $x+1=0$'' is true.  It is also true of the
one-element set, of a free group, and of the field of rational functions.  All of
them contain a solution; none of them is the integers.  So the sentence
\emph{there was no solution, therefore we added one} does not pick out an object.
It picks out a crowd.

This paper is about the missing part of that sentence.  We wanted a test that a
reader could apply to someone else's claim --- and, more usefully, to their own
--- and get an answer that does not depend on how the object was built.  We built
the test, checked its mathematical core in a proof assistant, and then pointed it
at our own programme.  It rejected one of our constructions, refused to certify
another, and broke on a third.  We report all three.

\paragraph{The idea in one sentence.}
An object is \emph{forced} when the demand cannot be met on the old carrier, and
the extension that meets it is the \emph{smallest possible} in a precise sense ---
smallest not in size, but in the sense that it maps into every other extension,
uniquely.

The word ``smallest'' needs care, because it is where most claims of necessity
quietly cheat.  Take the demand ``$x+1=0$ must become solvable''.  If we are
allowed to keep working among structures where addition is only a semigroup
operation, we can add a formal symbol $\varepsilon$ with $\varepsilon+1=0$ and
stop.  The result solves the equation and is not the integers: its addition is
not a group operation.  If instead we first agree to work among commutative rings,
the same demand produces the integers, and produces nothing else.  The difference
is not in the demand.  It is in the company the demand is allowed to keep.

So the test has to fix the company first.  This is the single most useful lesson
we took from the exercise, and it is easy to state: \emph{changing the ambient
theory changes the object, and the change of theory is the real content of the
step.}

\subsection*{Background: the same object, built many ways}

The number systems are the oldest worked example of this problem, and they are
worked precisely because the answer is not in the construction.  The integers can
be built as pairs of naturals modulo a difference, or as the group completion of
the additive monoid; the rationals as ordered pairs of integers, or as the
fraction field of the integers; the reals as Dedekind cuts (1872) or as
equivalence classes of Cauchy sequences; the complex numbers as pairs of reals
(Hamilton, 1837) or as the quotient ring $\mathbb{R}[x]/(x^2+1)$.  Every pair of
these constructions gives isomorphic objects.  Landau's \emph{Grundlagen der
Analysis} (1930) builds the tower from Peano's axioms, with the standard
uniqueness statements; Ebbinghaus and co-authors' \emph{Numbers} (1991)
surveys several constructions of each system together with their equivalence,
and is, as far as we can tell, the book a mathematician reaches for when they
want the whole tower in one place \cite{landau,dedekind,ebb}.

That literature settles \emph{what} the objects are.  It does not settle the
sentence that mathematicians actually use when they introduce them, which is
causal rather than structural: \emph{we could not subtract, so we made the
integers}.  The causal sentence is the interesting one, because it is the one
that gets abused.  A program that adds a formal square root of seven can say the
same sentence, word for word, and it is wrong.

The philosophical version of the difficulty is older than the formal one.
Benacerraf (1965) observed that if numbers are sets, then the natural number 4
could be the von Neumann 4 or the Zermelo 4, and nothing in arithmetic tells us
which; so either the identification is arbitrary or numbers are not objects at
all \cite{benacerraf1965}.  Structuralism replies that numbers are positions in a
structure rather than objects with an inner nature \cite{shapiro1997}.  We take
no position on that debate.  We take from it a design constraint: whatever test
we propose must give the same answer for every construction of the same object.
That constraint is what the second question in this paper enforces, and it is why
the answer to the test is a statement about a universal property rather than
about a set.

\subsection*{What has already been said about necessity}

Mathematicians do have a standard name for \emph{had to exist}: the object is
\emph{initial} (or \emph{free}, or \emph{universal}) in the relevant category.
Mac Lane's textbook states and uses the fact that an initial object is determined
up to a unique isomorphism \cite{MacLane1971}.  The specific instances are
textbook too: the integers are the group completion of the additive monoid,
the rationals are the fraction field, the complex numbers are the ring generated
by a root of $x^2+1$.  Model theory has a parallel notion: when a model completion exists it is unique
up to equivalence, and the comparison embeddings are elementary and unique over
the base model \cite{EklofSabbagh1971,Hodges1993}.

So the ingredients are not new, and we say so in the paper rather than letting a
referee say it.  As far as we can tell from the literature cited here, what is
missing is a \emph{test}: a short
list of questions that a reader can put to a claim of the form ``this object had
to exist'', including claims made by a program, together with an honest account
of which questions can be answered mechanically and which cannot.  The textbook
statement is about a fixed category.  A test has to survive a change of category,
because that is exactly where the abuses live.

\subsection*{The question, stated plainly}

\begin{table}[t]
\caption{The problem this paper addresses, in the form we would want a referee to
check.  The last row distinguishes three levels of evidence rather than claiming
one: checked, argued, asserted.}
\label{tab:problem}
\begin{ruledtabular}
\begin{tabular}{p{0.14\linewidth}p{0.74\linewidth}}
Item & Statement \\
\hline
Given & A structure $S$ (objects with operations) and a demand $D$ that $S$
cannot meet. \\
Asked & Is there an object that $D$ \emph{forces} --- one that every acceptable
answer must be, up to a unique relabelling? \\
Obstacle 1 & $D$ may already be met on the same objects, in which case nothing is
forced. \\
Obstacle 2 & Many extensions may meet $D$, including trivial and chaotic ones;
without a smallest one, ``the'' answer is not defined. \\
Obstacle 3 & The answer may be an existing object under a new name. \\
Our answer & Three questions, one per obstacle: same-carrier solvability;
existence of a smallest extension, unique up to relabelling; novelty against an
explicit catalogue. \\
Delivered & Three levels of evidence, kept apart: machine-checked (three tower
steps, the four failures, the two-score separation); argued by hand (the counting
refusal); asserted (the minimality of $\mathbb{R}$).  Question 3 is specified but
not implemented, and our first implementation was refuted by our own examples. \\
\end{tabular}
\end{ruledtabular}
\end{table}

With that, the question of the paper is narrow enough to answer.  Given a
structure and a demand, is there an object that the demand forces, and can we
tell --- mechanically, in a way a sceptical reader can repeat --- when the answer
is yes?

\subsection*{What we found, in advance}

Four results, stated here so that a reader who stops after this page still knows
what happened.

\begin{itemize}
\item \emph{Three steps of the tower pass.}  The integers, the rationals and
      $(\mathbb{C},i)$ are each the smallest extension of the previous system
      once the ambient theory is fixed, and each is determined up to a unique
      relabelling.  These are classical facts; the contribution is that one
      statement covers all three.
\item \emph{One step cannot be certified.}  The reals are not proved to be the
      smallest complete ordered field over the rationals.  We prove the failure
      that motivates the step and the cut property that describes the reals, and
      we register the missing statement as a named obligation that nothing
      depends on.
\item \emph{One candidate is refused.}  A counting invariant proposed inside our
      own programme collapses to a difference of prime-factor counts, hence is
      already definable on the same objects.  The test refuses it, by hand.
\item \emph{The third question failed.}  Our implementation of ``is it new?''
      compared the number of elements, which makes the integers, the rationals
      and the complex numbers renamings of one another; the gate then refuses
      everything, including the objects the rest of the paper certifies.
\end{itemize}

\section{What others have done}
\label{sec:related}

Four bodies of work bear on this question, and the boundary between them and this
paper is worth drawing explicitly.

\paragraph{Building the numbers, and proving the result unique.}
The constructions are nineteenth- to mid-twentieth-century mathematics:
Dedekind's cuts, Cantor's sequences, Hamilton's pairs, and the algebraic
constructions that turn a semiring into a ring and a ring into a field.  The
systematic account, with uniqueness proofs, is \cite{ebb}; the axiomatic account
from Peano's axioms is \cite{landau}.  This literature is the ground truth for
every example in this paper, and we claim nothing new about the objects
themselves.

\paragraph{Universal properties and the uniqueness of the comparison.}
The general statement --- an initial object is determined up to a unique
isomorphism --- is in \cite{MacLane1971}, and it is used constantly: free groups,
localisation, polynomial rings, quotients.  Model theory supplies the same shape
in a different language, through model completions and the uniqueness of the
comparison map \cite{EklofSabbagh1971,Hodges1993}.  Our main theorem is a
parameterised form of this statement: the ambient theory and the refinement are
arguments, so the same four-line argument covers the whole tower.  We are explicit
in the text that this is packaging, not new mathematics.

\paragraph{Formalisation.}
The reals have been built in proof assistants before --- in HOL Light
\cite{Harrison1998}, and earlier in HOL88 --- and the algebraic hierarchy used
here is Mathlib's \cite{mathlib}, in the Lean theorem prover \cite{lean}; this
development uses Lean~4.  Practical advice about the instance
parameters that make such a hierarchy usable is in \cite{Baanen2022}.  What a
proof assistant adds to a paper of this kind is not depth but \emph{auditability}:
the definitions cannot silently mean something else, and a claim that a proof is
four lines can be checked as a fact about the file rather than taken on trust.
In our experience this is worth the cost mainly because it makes mistakes visible
--- two of ours are in this paper.

\paragraph{Machine discovery.}
Programs that propose mathematics have a long history, from Lenat's
\textsc{AM} \cite{lenat1976} to recent systems that use search or learning to
generate conjectures and constructions \cite{davies2021,raayoni2021,romera2024}.
Their pattern is consistent: generate candidates, then decide which ones are
worth keeping.  The decision is made by a human, by a heuristic of interestingness
(as in \textsc{AM}), or by a checker that tests truth or performance rather than
novelty.  This paper is about the decision step.  It does
not generate anything, and it should be read as complementary to those systems: a
generator supplies candidates, a criterion decides which of them may be called
forced.

\paragraph{Philosophy: discovery, invention, and what numbers could not be.}
The debate this paper's question belongs to is Benacerraf's
\cite{benacerraf1965}, with the structuralist response in \cite{shapiro1997}.
We do not try to settle it.  We take one thing from it: if the answer is to be
independent of the construction, the test must be invariant under change of
construction.  That is a formal requirement, and it is what Question~2 enforces.
The philosophical reading of ``forced'' remains a reading; what we can check is
that the object is determined.

\paragraph{Where this paper sits.}
Not in the construction of the objects, not in the uniqueness theorem, and not in
the formalisation.  It sits in the combination: a short test, stated so that it
applies to any theory, applied honestly to the author's own programme, with the
negative results reported at the same length as the positive ones.  The closest
neighbour is the category-theoretic statement of initiality; the difference is
that we ask what happens when the initiality claim is \emph{false or unproved},
which is the situation a working mathematician is usually in.


\section{A test with three questions}
\label{sec:test}

We fix, before asking anything, three pieces of data: a theory (what the new
object may carry), a way of reading the old structure inside the new one, and a
demand (a constant to add, a property to impose, or both).  With those fixed, the
test is three questions.

\paragraph{Question 1: is the demand really unmet?}
The question is about the elements, not about the solution.  The wording is doing
work.  Both the natural numbers and the integers are countable, so any ring
structure on the integers can be carried over to the natural numbers by a
bijection.  The statement ``the natural numbers admit no ring structure'' is
therefore \emph{false}, and an early version of our own formalisation asserted it.
The true statement keeps the old structure fixed: no ring structure on
$\mathbb{N}$ whose addition and multiplication extend the usual ones solves
$x+1=0$.  This episode is worth recording because the test caught its author.

\paragraph{Question 2: is the extension the smallest one?}
Among all extensions that meet the demand, is there one that maps into every
other, and does exactly one such map exist?  When the answer is yes we call the
extension \emph{minimal}.  The name is traditional; the property is not
smallness but universality, and universality is what makes the object determined.
The next section states what follows.

\paragraph{Question 3: is it new?}
New relative to what is already known --- not relative to the system's own state,
which can always be arranged to report novelty.  In practice this means an
explicit catalogue of known objects, supplied from outside.  Section~\ref{sec:gate}
reports what happened when we tried to implement this question; the short version
is that our implementation was wrong in a way that only the examples revealed.

\begin{table}[t]
\caption{The three questions, what each one rules out, and what we actually
delivered.  ``Answered'' means a machine-checked statement about the examples in
this paper; ``specified'' means we wrote down what the answer would have to be but
did not implement it.}
\label{tab:three}
\begin{ruledtabular}
\begin{tabular}{p{0.06\linewidth}p{0.34\linewidth}p{0.3\linewidth}p{0.16\linewidth}}
\# & Question & Rules out & Status here \\
\hline
1 & Can the old carrier support the demand? & changing the carrier and calling it a solution & answered \\
2 & Is there a smallest extension, unique up to relabelling? & the trivial extension, the chaotic extension, and every construction that depends on choices & answered \\
3 & Is it new, given a catalogue of known objects? & renaming an existing object & specified, not implemented; our first implementation failed (\S\ref{sec:gate}) \\
\end{tabular}
\end{ruledtabular}
\end{table}

\subsection*{What ``forced'' means here, and what it does not}

The word carries at least four different claims, and a paper that does not say
which one it means can be attacked from any direction.  In order of strength:

\begin{itemize}
\item \emph{Metaphysical}: the object could not have failed to exist.  We do not
      touch this.
\item \emph{Conceptual}: the object is already contained in the concept we
      started from.  We do not touch this either; our demands are data, not
      concepts.
\item \emph{Practical}: no working alternative exists, given the state of the
      field.  That is a claim about a community, not about mathematics.
\item \emph{Determined}: every acceptable answer is the same object up to a
      unique relabelling.  This is the one we test.
\end{itemize}

The fourth is weaker than the first three, and it is the only one that can be
settled by reading a proof.  It is also the one Benacerraf's argument leaves
standing: if numbers are positions in a structure, then ``the same position'' is
exactly determination up to relabelling, and nothing stronger is on offer.  We
therefore do not claim to have shown that the integers had to exist.  We claim to
have shown that, given the demand and the company it keeps, nothing else could
have answered it.

That is not a hedge; it changes what counts as a refutation.  A reader who wants
metaphysical necessity will find the test silent, and should.  A reader who wants
to know whether a program's ``discovery'' is a renaming will find it decisive.

\section{Why the second question settles the object}
\label{sec:theorem}

Question~2 has a clean answer, and the answer is the reason the test is worth
running.  If an extension is smallest in the sense above, then:

\begin{enumerate}
\item it maps into every other extension, and the map is unique;
\item any two smallest extensions are the same up to a relabelling of their
      elements, and that relabelling is itself unique;
\item the smallest extension has no symmetries of its own beyond the identity.
\end{enumerate}

None of this mentions rings, order, or completeness.  The three statements use
four properties of the surrounding category --- an identity map, a composition,
an underlying-function map, and the fact that maps with the same values are equal
--- and nothing else.  That is why one argument covers the integers, the
rationals and the complex numbers at once, and why the argument can be checked
once rather than three times.

We should be plain about how much this is worth.  Once you \emph{define}
``smallest'' as ``maps into everything, uniquely'', the first and third
statements are restatements of the definition, and the proof is four lines.  The
content is not the proof.  The content is that this is the right definition to
have chosen: it is the only one on the table that does not depend on how the
object is constructed.  Two people who build the complex numbers differently --- as
pairs of reals, as a quotient ring, as a matrix algebra --- get the same answer to
Question~2, and their constructions are related by a relabelling that is itself
unique.  That is what it means for the object, rather than the construction, to be
determined.

\paragraph{A worked example, and a warning.}
For the step from the reals to the complex numbers, the demand is ``add a constant
$j$ with $j\cdot j=-1$''.  The extension $(\mathbb{C},i)$ is smallest, and it has no
symmetry beyond the identity.  Now drop the constant from the demand and ask only
that \emph{some} square root of $-1$ exists.  The same field is still an
extension, but it is no longer smallest: the comparison is no longer unique.
Identity and complex conjugation are two distinct maps over $\mathbb{R}$, and both
are allowed; pinning $i$ is what removes the second one.  Finding a solution and
pinning it down are different achievements, and only the second produces an
object.  A test that accepts the unpinned version will accept a field with a
symmetry group, which is exactly the ambiguity it was supposed to remove.

\paragraph{One thing the argument does not settle.}
The demand ``the order is complete'' behaves differently from ``$j^2=-1$'', and it
is worth saying why in words rather than symbols.  An equation like $j^2=-1$ is
preserved by every ring map: if $j$ satisfies it, so does the image of $j$.  A
condition like ``every bounded non-empty set has a least upper bound'' is not
preserved by ring maps, and it has no term model --- there is no expression you can
write down, quotient by a relation, and be done.  What follows is that the
smallest extension cannot be \emph{computed} the way the complex numbers can.  It
does not follow that no smallest extension exists.  Our first draft drew the
stronger conclusion; the referee was right and we withdrew it
(Sec.~\ref{sec:real}).

\section{What the test found}
\label{sec:found}

We ran the test on four steps of the number tower, and on two further pieces of
our own programme: an invariant built from prime factorisations, and a coordinate
built from reversible dynamics.  Table~\ref{tab:results} gives the summary; the
two rows marked ``not applicable'' are those pieces, and the reason is that
neither is a demand on a structure --- one is a definition, the other a coordinate
system --- so Questions 1 and 2 have nothing to bite on.

\begin{table}[t]
\caption{What passed, what did not, and what could not be judged.  ``Same
carrier'' is Question~1, ``smallest'' is Question~2, ``new'' is Question~3.
``Proved'' means machine-checked modulo the standard-maps condition of
Sec.~\ref{sec:limits}; ``asserted'' means written down as an obligation and not
proved; a dash means the question does not apply or was not used.  No entry in the
``New'' column was produced by an implementation: the novelty numbers printed by
the shipped trace were entered by hand.}
\label{tab:results}
\begin{ruledtabular}
\begin{tabular}{lccc}
Case & Same carrier & Smallest & New \\
\hline
$\mathbb{N}\to\mathbb{Z}$ & proved & proved & not used \\
$\mathbb{Z}\to\mathbb{Q}$ & proved & proved & not used \\
$\mathbb{Q}\to\mathbb{R}$ & proved & \emph{asserted} & not used \\
$\mathbb{R}\to\mathbb{C}$ & proved & proved & not used \\
powers vs roots & not applicable & not applicable & not applicable \\
counting invariant & \emph{fails} (refused by hand) & --- & --- \\
vibration coordinate & not applicable & not applicable & not applicable \\
\end{tabular}
\end{ruledtabular}
\end{table}

\subsection{Three steps that pass}
\label{sec:tower}

The integers are the smallest extension of the naturals once we agree to work among
commutative rings; the rationals are the smallest extension of the integers once we
agree to work among fields; and $(\mathbb{C},i)$ is the smallest extension of the
reals once we pin the square root of $-1$.  The underlying facts are classical (compare \cite{ebb}); the formulation as a
smallest extension, with the old structure pinned, is our packaging.  What the
test adds is that the three are instances of one statement, and that in each case
the answer does not depend on the construction.  For the integers, the
side condition that the extension really contains the old system also gives
characteristic zero for free, which is a pleasant sign that the definition is
carrying its weight.

\subsection{One step we could not certify}
\label{sec:real}

For the step from the rationals to the reals we prove the failure --- the rationals
are not complete --- and we prove that every real is generated by a cut of
rationals below it.  We do \emph{not} prove that the reals are the smallest
complete ordered field over the rationals.  That statement is written down as a
named obligation in the formal development, and nothing in the development uses
it.

We think the obligation is true and provable; a page of standard arguments about
complete ordered fields should do it.  We did not do it.  The consequence is that
the fourth step is certified less strongly than the first three, and we prefer to
say so in the text rather than to let a reader discover it in the appendix.  The
asymmetry is about proof debt and about what can be computed, not about a
metaphysical difference between adjoining a root and completing an order.

\subsection{One candidate we refused}
\label{sec:counting}

Our programme also proposed an invariant of integers, meant to capture the
``tension'' between two numbers: count the primes that divide one but not the
other, subtract the count in the opposite direction, and call the result
$\tau(a,b)$.  The proposal was that the surface $\tau=0$ organises arithmetic, and
that $\tau$ is a new object in its own right.

The axiomatic part worked out cleanly.  Under four natural assumptions, the
invariant is unique up to sign and equals a difference of counts,
$\tau(a,b)=\omega(a)-\omega(b)$, where $\omega(n)$ is the number of distinct prime
factors of $n$.  That is the whole content: $\tau$ knows how many prime factors
each side has and never which ones.  Since $\omega$ is a century-old function and
the formula is definable on the same carrier, the test returns \emph{refused} at
Question~1.  There is no new object here; there is a theorem about $\omega$, and a
clear statement of why the construction cannot say more.

Two honest qualifications.  The refusal is argued by hand: the counting part of
the development contains no demand and no catalogue, so ``Question~1 returns
refused'' is a mathematical observation, not a machine-checked verdict.  And the
refusal does not go through Question~3; our first draft said the novelty score was
zero there, which does not follow and is not available, since no catalogue exists.

The episode is still useful.  It is the case where the test said no to us, and it
came with a reason: the invariant collapses to a count, and counts are already on
the same carrier.  A weaker test would have produced a paper claiming a discovery.

\section{The measurement that failed}
\label{sec:gate}

Question~3 is the one we could not deliver, and we describe the failure because it
is instructive.

The design had two scores.  One was a weighted average of five qualities ---
coverage, relatedness, closure gain, minimality, novelty --- meant to \emph{rank}
candidates.  The other was a product of four --- novelty, consistency,
generativity, compression --- meant to \emph{admit} them.  The distinction is real:
a weighted average gives a high mark to a candidate whose novelty is zero but whose
other qualities are excellent, which is precisely the $\sqrt 7$ story; a product
gives it zero.  We checked the arithmetic of this distinction in the proof
assistant: there is a score with novelty zero and generativity $4/5$, and a
snapshot with the same entries whose product is zero
(\lean{weighted_sum_admits_renaming}, \lean{product_rejects_renaming}).  We stress
what this does and does not show.  It shows that the two formulas are not
interchangeable.  It does not model renaming at all: no catalogue, no candidate, no
map between objects appears in those two statements.

Modelling renaming is where the implementation failed.  To decide whether a
candidate was new, the code compared the \emph{sizes} of carriers: if a candidate's
carrier had the same cardinality as a known object's carrier, it was called a
renaming.  The integers, the rationals and the complex numbers are all countably
infinite.  One countable object in the catalogue therefore makes all three
renamings of it, and the gate assigns all three a novelty of zero.  The gate
refuses exactly the objects the rest of the paper certifies.

We did not patch this.  Comparing sizes can only produce false refusals, never
false admissions, so the risk runs in the opposite direction from the one our first
draft described.  A correct comparison is between structures, not carriers, which
requires the catalogue to carry its own maps; that is a specification, and we have
not implemented it.  Consequently no admission decision in this paper uses
Question~3, and the executable trace shipped with the code prints novelty values
that were entered by hand.  We put this in the table caption rather than in a
footnote, and the caption says it.

\section{What the test cannot do}
\label{sec:limits}

\paragraph{It assumes the standard maps.}
The four properties that the main argument uses do not say that the maps between
structures are the structure-preserving ones.  A degenerate choice --- one map
between any two structures --- satisfies all four, and in that setting every object
is both smallest and largest, so the test says nothing about anything.  Our
examples use the standard maps by construction, but the assumption is not provable
inside the framework, only checkable in each instance.  We mark every statement that
needs it as ``modulo C0''.  This is the price of stating the criterion for an
arbitrary theory instead of for rings.

\paragraph{It does not generate the demands.}
The hard parts of a discovery engine --- deciding whether a demand is satisfiable,
proposing a candidate, deciding minimality, deciding consistency --- are inputs
here, not code.  We prove what surrounds them: that a solvable demand never needs
a new carrier, that a contradiction is never expanded, that the old system embeds
in the new one, and that two discoveries of the same demand agree up to a unique
relabelling.  The engine is a judge, not an author.

\paragraph{It does not decide its own questions.}
Each question is a proposition about a fixed triple of data.  Deciding it is
per-instance work, and for Question~1 it is the same work as solving equations in
a structure.  Nothing here is a decision procedure, and we do not claim that the
three questions can be answered mechanically in general.

\paragraph{It has only refused one of our candidates.}
The other case, the reals, is an uncertified obligation rather than a refusal.  A
test that has said no once is a start, not a track record.  The strongest evidence
we could offer would be a refusal of a candidate proposed by someone else, and
believed by them; we do not have that yet.

\paragraph{The test is relative to the data.}
Every verdict in this paper is relative to three choices: the ambient theory, the
refinement, and --- for Question 3 --- the catalogue.  Change any of them and the
verdict can change.  This is not a defect to be patched; it is what makes the
verdict checkable.  But it does bound the claim: we have not shown that an object
is forced unconditionally, only that it is forced relative to a stated triple.  A
reader who wants the unconditional word will not find it here, and the honest
response is that we do not know what the unconditional word would mean.

\section{Discussion}
\label{sec:discussion}

The practical lesson is about the order in which to explain an object.  The usual
order is: here is the object, here are the properties that make it necessary.  The
order that survived our test is the reverse.  First the failure that cannot be
repaired in place.  Then the company the repair is allowed to keep, because the
same failure repaired in different company gives different objects.  Then the
smallest repair, which is determined once the company is fixed.  Whatever is left
over after those three steps is not forced, and should be called a construction.

The second lesson is that a negative result is cheap and worth publishing.  Our
programme's counting invariant was, for a while, the centrepiece of a different
paper.  Running it through Question~1 took an afternoon and produced a theorem
about $\omega$ instead.  That is a better outcome than a paper claiming a
discovery, and it cost less.

The third lesson is about machine-checking.  The proof assistant did not settle any
philosophical question, and it did not make the mathematics deep.  It did three
useful things: it forced the definition of ``gap'' to be stated in a way that is not
false; it made the four-line nature of the main argument visible, so we could not
pretend it was more; and it turned one of our claims about two scoring formulas
into a theorem that a reader can check in a minute.  It also made our mistakes
visible --- the ring-structure error, the carrier-counting error --- which is the
main reason to use it in a paper of this kind.

\paragraph{What would change our mind.}
Three things.  A genuine discovery that passes Questions~1 and~2 and is refused by a
repaired version of Question~3 would show that the third question is wrong, not
merely unimplemented.  A demand whose extensions do not form a category would show
that the second question is not available in general.  And a proof that the reals
are not the smallest complete ordered field over the rationals would make our
fourth row false rather than unproved.  We expect the third to go the other way:
the obligation looks provable, and we would be glad to withdraw the asymmetry.

\paragraph{How this relates to other work.}
The uniqueness half of Question~2 is category theory folklore \cite{MacLane1971},
and its model-theoretic cousin --- the uniqueness of the comparison supplied by a
model completion --- is in \cite{EklofSabbagh1971}.  What is added here is the
packaging: the same statement for an arbitrary theory, with the demand and the
ambient theory as parameters, and the two negative verdicts that follow from it.
Automated conjecturing systems \cite{lenat1976,davies2021,raayoni2021,romera2024}
generate statements and are judged by whether the statements are true and useful.
They supply candidates; this paper is about the other end, deciding which candidates
may be called forced.

\section{Methods}
\label{sec:methods}

\paragraph{Formal setting.}
A theory assigns to each type the type of structures on it; a model is a carrier
together with a structure.  A framework supplies maps between models, an identity,
a composition, an underlying-function map, and extensionality (maps with equal
values are equal).  Associativity and the unit laws for composition are derived
from these, not assumed.  A demand supplies, for each model, a type of witnesses,
a predicate on witnesses, and a transport of witnesses along maps that is
functorial.  An extension of a model is a target model, an injective map from the
old carrier, and a witness satisfying the predicate.  It is \emph{minimal} when it
maps into every other extension, uniquely.  The three conclusions of
Sec.~\ref{sec:theorem} are Theorem~1 of the development
(\lean{IsMinimal.hom_unique}, \lean{IsMinimal.carrierEquiv}, \lean{IsMinimal.rigid}).

\paragraph{Trust base.}
The development is Lean~4 \cite{lean} with Mathlib \cite{mathlib}, with no \texttt{sorry}.  There is one named
axiom, the initiality of the reals, and no theorem consumes it.  Separately,
\lean{native_decide} occurs in two files of the counting layer; it introduces
\lean{Lean.ofReduceBool} into the trusted fragment, and it supports finite numerical
checks that this paper does not report.  The statements this paper does report from
that layer are proved by rewriting and arithmetic, not by evaluation; our earlier
draft stated this the other way round and the correction is recorded in the
repository.

\paragraph{What is checked, and what is only argued.}
Checked in the development (each has a named declaration; the names are collected
in the audit file): the main theorem; the four failures; the three minimality
results; the characteristic-zero side condition; the exclusion lemmas for two of
the three pairs of verdicts; the engine's four positive properties; the torsion
decomposition of the power map; the collapse of the counting invariant to a count;
the two-score separation.  Argued but not checked: the preservation of equations versus closure
conditions; the refusal of the counting invariant; the comparison between a
candidate and the old world in the gap branch.  Assumed: the standard-maps
condition, and the initiality of the reals.

\paragraph{Data and code.}
The Lean development and the seven-round self-audit of this manuscript are in the
repository, tagged \lean{paper-forced-objects-v1.1} (commit \lean{720ae3b}), and the
build pins Mathlib in
\lean{lakefile.pinned.toml}; every declaration named in the text can be located by its identifier,
and the audit records which claims were withdrawn or downgraded, and why.  Three
limits on that availability are stated rather than glossed.  There is no separate
executable model for the counting or vibration layers: those are Lean
definitions, and the trace in \lean{Nods/Chain.lean} is a Lean file whose scores
were entered by hand, not a program that computes them.  The repository also
contains unrelated work (a simulation study) that is not part of this paper, and
we have not pinned the Mathlib revision, so we do not advertise a one-command
buildable artifact.

\begin{acknowledgments}
The development was carried out in Lean~4 with Mathlib, and thanks are due to the
Mathlib community.  The manuscript was put through seven rounds of adversarial self-audit, run as
independent referee lines by the author's tooling; the record, including the
withdrawal of two factual claims and the demotion of the third question, is kept
with the source.
\end{acknowledgments}

\section*{Competing interests}
The author declares no competing interests.

\bibliography{nods}

\end{document}
